% TODO(team): Explain exposure units and the distinction between grace and cancellation.
\section{E2: repair timing}
\begin{frame}{Grace trades early action against unnecessary work}
\begin{itemize}
\item Shared repair policy S; grace of \textbf{1 / 3 / 8 seconds}.
\item Brief absence: original Peer returns after \textbf{4 seconds}.
\item Persistent absence: original Peer stays offline.
\item Measure repair payload and time spent with one reachable copy.
\end{itemize}
\takeaway{Grace controls when repair starts, not whether a started task is cancelled.}
\note{An unavailable Peer is not necessarily permanently lost. Grace can avoid copying data that will soon return, but keeps fewer reachable copies for longer.}
\end{frame}
\begin{frame}{Waiting avoids copies, but extends exposure}
\small
\begin{tabularx}{\textwidth}{lYYY}
\toprule Grace & Brief: repair MiB & Persistent: repair MiB & Persistent exposure MiB$\cdot$s \\\midrule
1 s & 0.125 & 0.750 & 17.42\\
3 s & 0 & 0.750 & 18.24\\
8 s & 0 & 0.750 & 20.20\\\bottomrule
\end{tabularx}
\vspace{4mm}\normalsize
All 54 brief-outage viewer sessions completed.
\takeaway{Persistent exposure increased about 16\% from 1 s to 8 s grace.}
\note{One-second grace copied one 128 KiB segment per brief-outage run. Longer grace avoided that work. Full restoration is not strictly monotonic in every context because source contention and coordinator ticks also matter. Source: compact-study-results.md, E2.}
\end{frame}
